\documentclass[11pt]{article}
\usepackage[T1]{fontenc}
\usepackage[utf8]{inputenc}
\usepackage{lmodern,amsmath,amssymb,amsthm,mathtools}
\usepackage[margin=1in]{geometry}
\usepackage{booktabs,microtype,listings,enumitem,fancyhdr}
\usepackage[hidelinks,breaklinks]{hyperref}
\usepackage{xurl}
\hypersetup{pdftitle={Crossed uniqueness and the concise iPRA frontier},pdfauthor={OpenAI, prepared for Vladimir Reshetnikov},pdfsubject={A mathematical primitive-basis proposal and a measured byte refinement}}
\setlength{\emergencystretch}{3em}
\setlength{\parskip}{3pt}
\setlist{itemsep=2pt,topsep=4pt}
\pagestyle{fancy}\fancyhf{}
\fancyhead[L]{\small Crossed uniqueness for iPRA}
\fancyhead[R]{\small 12 November 2026}
\fancyfoot[C]{\thepage}\setlength{\headheight}{14pt}
\newtheorem{theorem}{Theorem}[section]
\newtheorem{lemma}[theorem]{Lemma}
\newtheorem{proposition}[theorem]{Proposition}
\newtheorem{corollary}[theorem]{Corollary}
\newcommand{\Q}{\mathsf Q}\newcommand{\D}{\mathsf D}
\newcommand{\U}{\mathsf U}\newcommand{\nf}{\operatorname{nf}}
\newcommand{\one}{S0}
\newcommand{\file}[1]{\nolinkurl{#1}}
\lstset{basicstyle=\ttfamily\footnotesize,breaklines=true,breakatwhitespace=false,
 columns=fullflexible,keepspaces=true,showstringspaces=false,frame=single}
\title{\textbf{Crossed uniqueness and the concise iPRA frontier}\\[5pt]
\large A uniform arithmetic rule and a measured 502-byte refinement}
\author{Prepared for Vladimir Reshetnikov\\
\small Research report with reproducible artifacts}
\date{12 November 2026}
\begin{document}
\maketitle
\begin{abstract}
On computationally normalized scalar-iterator terms, an arbitrary-seed
crossed uniqueness rule derives all equality transport as well as equation
induction. With the native case logic it gives a six-rule presentation
bi-interpretable with full first-order iPRA. The new elementary argument
derives symmetry, congruence and transitivity before induction. The
subsequent scalar uniqueness-to-induction bridge is inherited and
explicitly attributed.

Separately, a composition of existing implementation refinements reduces
the smallest audited compact byte checker from 506 to 502 normalized bytes.
This program has 269 significant tokens and 301 AST nodes. It checks the
existing guarded presentation, not the new uniqueness calculus. Its strict
input-validating companion has 1172 bytes.

The requested improvement below the complete current complexity envelope
has \emph{not} been established. The prior envelope is 253 tokens,
275 AST nodes and 506 bytes, attained by different programs. The new
mathematical checker is larger. Neither absolute minimality nor a new
Lean or Rocq formalization is claimed.
\end{abstract}
\tableofcontents
\newpage

% BEGIN INCLUDED sections/frontier.tex
\section{Result, comparison target, and claim boundaries}
The source snapshot is ProveIt commit
\begin{center}\small\texttt{6ffc49f55d355b548f7a8e1530570aee5de086b3}.\end{center}
The overview alone does not identify the actual comparison target.
All three incoming archives were inspected. A frozen source census
contains 1380 distinct Python Git blobs at 5825 paths; all 90 current
Research package tree identities match it. Its measurements and the
classification of 106 smaller utility modules were reproduced.

The incoming \file{ipra_smaller_envelope.zip} supplies the token and AST
records. The source
\file{deterministic_introduction/implementation/code/compact_optional.py}
supplies the byte record, Git blob
\texttt{a0af7afb236b530eb30c63c5ab788151024271b7}.
This 506-byte incumbent was not identified as the byte record in the
overview table. The full census and comparison sources are included.

\begin{center}
\begin{tabular}{lrrr}
\toprule
Complete module & Tokens & AST nodes & Bytes\\
\midrule
Incoming token profile & \textbf{253} & 280 & 599\\
Incoming AST profile & 257 & \textbf{275} & 607\\
Audited byte control & 267 & 303 & \textbf{506}\\
New compact composition & 269 & 301 & \textbf{502}\\
Strict validating control & 462 & 524 & 1183\\
New strict companion & 463 & 519 & 1172\\
New uniqueness reference checker & 1909 & 2087 & 5618\\
\bottomrule
\end{tabular}
\end{center}

The census convention uses CPython 3.12.14: remove comments and genuine
module/function/class docstrings, normalize with \file{ast.unparse} plus
one final newline, count significant tokens and \file{ast.walk} nodes,
and count UTF-8 bytes. The archive also supplies an actual normalized
executable of the new reference checker. Its documented source, with
docstrings retained under literal \file{ast.unparse}, measures
1914 tokens, 2097 AST nodes and 6686 bytes. No trusted code is omitted.

The envelope $(253,275,506)$ is empirical, not a lower-bound theorem for
all possible axiomatizations, and it is not the size of one program.
The compact composition lowers only the byte coordinate to 502, a
four-byte or approximately $0.79\%$ improvement.

\subsection{Input contracts}
The compact modules assume a finite integer/list certificate and a
canonical closed guarded Eq/D public goal. They do not themselves
validate that entire input promise. Their strict companions count
alphabet and public-grammar validation within the complete module.
The uniqueness checker has a different explicit syntax, validates open
sequents as well as closed theorems, and counts normalization,
substitution, scope checks, all rules and JSON input/output.

The mathematical contribution is the fusion of equality with uniqueness
through an arbitrary initial equation and a crossed conclusion.
Uniqueness as an alternative to induction is classical
\cite{goodstein,tait}; the scalar completion is already in ProveIt
\cite{ur}. The byte result composes three already proposed refinements:
delayed elimination, eager tuple selection, and the shorter mode selector.
It implements the existing calculus, not the new mathematical basis.

No complete implementation of the new basis below the current envelope
is established. In particular, six rule names do not beat existing
three-family clause calculi by themselves. The report supplies a
mathematical primitive-basis proposal and a smaller byte checker, while
leaving the requested full complexity breakthrough open.


% END INCLUDED


% BEGIN INCLUDED sections/calculus.tex
\section{A complete six-rule presentation}\label{sec:system}
\subsection{Syntax and computation}
There is one numerical sort with raw terms
\[
 t::=x\mid0\mid S(t)\mid I(n,a;u.h).
\]
Only $u$ in $h$ is bound. Updates are fixed finite numerical templates
with passive parameters; there are no function variables or universal
evaluators for primitive-recursive descriptions.

Represent terms in full normal form for
\[
 I(0,a;u.h)\longrightarrow a,\qquad
 I(Sn,a;u.h)\longrightarrow h[I(n,a;u.h)/u].
 \tag{Comp}\label{eq:comp}
\]
Every substitution is capture avoiding and followed by normalization.
For $C(v,a,b):=I(v,a;u.b)$ with unused $u$,
\[
 C(0,a,b)\equiv a,\quad C(Sv,a,b)\equiv b,\quad
 I(1,a;u.h)\equiv h[a/u].
\]
Definitional equality here uses only the two computation equations,
not arbitrary proved arithmetic equalities.

Normalization is counted as part of the trusted specification.
The repository supplies totality, canonical forms and
\[
 \nf(\nf(t)[\sigma])=\nf(t[\sigma]).                    \tag{NCoh}
\]
Its normalization-by-evaluation proof interprets a structurally smaller
body as a total operation on normal terms, traverses the finite
successor prefix of the normalized counter, and retains an iterator at
a neutral counter. Environment laws establish substitution coherence.
The reference checker uses equivalent direct reductions with capture
avoidance and memoization \cite{normalization,current}.

Formulas and their ordinary readings are
\[
 A::=(s=t)\mid\Q_yA\mid\D_{x,y}A,
\]
\[
 \Q_yB:=B(0)\to\forall y\,B(y),\qquad
 \D_{x,y}B:=\exists x\bigl(B(x,0)\to\forall y\,B(x,y)\bigr).
 \tag{Read}\label{eq:read}
\]
D binds distinct $x,y$. Contexts are finite reusable lists of formulas.
The native contradiction is $0=S0$. The readings do not silently insert
ordinary equality logic.

\subsection{All primitive rules}
All displayed expressions are normalized. The six families are
\[
 \frac{A\in\Gamma}{\Gamma\vdash A}\ {\rm Hyp}
 \qquad
 \frac{\ }{\Gamma\vdash t=t}\ {\rm Ref},
\]
\[
 \frac{\Gamma\vdash\Q_yB(y)\qquad\Gamma\vdash B(0)}
      {\Gamma\vdash B(t)}\ {\rm QE},
\]
\[
 \frac{\Gamma\vdash\D_{x,y}B(x,y)\qquad
       \Gamma,\Q_yB(b,y)\vdash A}
      {\Gamma\vdash A}\ {\rm DE},
\]
where $b$ is fresh for $\Gamma,A$ and the passive parameters of $B$;
\[
 \frac{\Gamma,B(t,0)\vdash B(t,Sz)}
      {\Gamma\vdash\D_{x,y}B(x,y)}\ {\rm DI_c},
\]
where $z$ is fresh for $\Gamma,t$ and the passive parameters of $B$; and
\[
 \boxed{\quad
 \frac{\Gamma\vdash f(0)=a\qquad
       \Gamma\vdash f(Sz)=h(f(z))}
      {\Gamma\vdash I(n,a;u.h(u))=f(n)}\ \U .
 \quad}                                                \tag{U}\label{eq:U}
\]
In U, $f(z),h(u)$ are fixed finite numerical templates. The eigenparameter
$z$ is fresh for $\Gamma,a,n,h$ and all passive parameters of $f$.
Thus $h$ has no active-counter argument. The recurrence premise has no
induction hypothesis. DI$_c$ always uses cases, even at equation matrices.

There are no other logical axioms, computation axioms, transport rules
or induction rules. Computation is exactly \eqref{eq:comp} in the syntax.

\begin{lemma}[Structural admissibility]
Weakening, renaming, ordinary-assumption cut and numerical substitution
followed by normalization are admissible.
\end{lemma}
\begin{proof}
Induct on finite derivations, freshening eigenparameters first.
Normalization coherence identifies substituted premises and conclusions.
Neither DI$_c$ nor U has a mode that can change under substitution.
Replace each use of a cut assumption by its proof. These are
metatheoretic inductions on finite syntax, not numerical induction.
\end{proof}

\section{Equality derived before induction}\label{sec:equality}
\subsection{Numerical cases without transport}
\begin{lemma}
The case logic, without U, derives
\[
 \frac{\Gamma\vdash A(0)\qquad\Gamma\vdash A(Sz)}
      {\Gamma\vdash A(n)}
 \quad\text{and}\quad
 \frac{\Gamma,A(0)\vdash A(Sz)}
      {\Gamma\vdash\Q_yA(y)},                          \tag{Cases}\label{eq:cases}
\]
with fresh $z$.
\end{lemma}
\begin{proof}
DI$_c$ with an unused coordinate $w$ introduces $\D_{w,y}A(y)$.
For the first rule weaken the successor proof by $A(0)$.
DE supplies $\Q_yA(y)$, and QE with the base gives $A(n)$,
independent of the opened witness. For the second rule DE exports the
opened Q assumption itself by Hyp.
\end{proof}

\subsection{Symmetry}
Given $s=t$, choose
\[
 f(z)=s,\qquad a=t,\qquad h(u)=u,\qquad n=0
\]
in U. Its recurrence is Ref. Its conclusion normalizes to $t=s$.

\subsection{Congruence in any numerical template}
For a template $H(u)$ and given $s=t$, choose
\[
 f(z)=I(z,s;u.H(u)),\qquad a=t,\qquad h=H,\qquad n=1.
\]
The base is $s=t$, while both sides of the recurrence normalize to
$H(I(z,s;H))$. Ref proves it. The conclusion is $H(t)=H(s)$;
symmetry gives ordinary congruence. Templates may duplicate or omit
the argument, or place it beneath additional iterator binders.

\subsection{Joining and transitivity}
\begin{lemma}[Common-source joining]
From $s=t$ and $s=u$ in one retained context, derive $t=u$ without
numerical induction.
\end{lemma}
\begin{proof}
Let $f(z)=C(z,s,u)$ and use the identity update. Its recurrence is
\[
 u=C(z,s,u).
\]
At zero this is $u=s$, by symmetry of $s=u$; at a successor it is
$u=u$, by Ref. Numerical cases prove the recurrence.
U with $a=t,n=1$ has base $s=t$ and conclusion $t=u$.
\end{proof}
From $s=t,t=u$, first derive $t=s$ and join at $t$ to get $s=u$.
Thus transitivity is derived. From $s=t,H(s)=r$, congruence supplies
$H(s)=H(t)$ and joining gives $r=H(t)$, the repository's one-term
crossed transport rule.

\subsection{Full formula transport}
\begin{theorem}[Leibniz transport]
For every native formula template $F(v)$,
\[
 \Gamma\vdash s=t,\quad\Gamma\vdash F(s)
       \quad\Longrightarrow\quad\Gamma\vdash F(t).
\]
\end{theorem}
\begin{proof}
Induct on $F$. At an equation use congruence on both sides, symmetry and
transitivity.

For $F(v)=\Q_yB(v,y)$, temporarily assume $B(t,0)$.
Reverse matrix transport gives $B(s,0)$; QE at $Sz$ gives $B(s,Sz)$;
forward matrix transport gives $B(t,Sz)$. Conditional Q introduction
\eqref{eq:cases} finishes.

For $F(v)=\D_{x,y}B(v,x,y)$, open at fresh $b$ and retain
$\Q_yB(s,b,y)$. Introduce the desired D formula at witness $b$.
From its base $B(t,b,0)$, reverse matrix transport gives $B(s,b,0)$,
QE gives $B(s,b,Sz)$, and forward transport gives $B(t,b,Sz)$.
DI$_c$ closes the desired formula and DE discharges $b$.
Each binder case uses only the strictly smaller matrix.
\end{proof}
This elementary fusion theorem assumes neither predecessor, numerical
equality tests, successor injectivity nor an induction-hypothesis rule.


% END INCLUDED


% BEGIN INCLUDED sections/bridge.tex
\section{The full logical and arithmetic bridge}\label{sec:bridge}
\subsection{All intuitionistic operations}\label{sec:choice-introduction}
The intended comparison is with full first-order iPRA. The following
macros are the inherited induction-independent case-logic construction
\cite{choice}; they now apply because equality was already recovered.

Replace Q inside macros by its interderivable D form with unused witness.
On Q-free formulas define
\[
 G_v(s=t):=(C(v,0,s)=C(v,0,t)),\qquad
 G_v(\D_{x,y}B):=\D_{x,y}G_vB.
\]
With capture avoidance, computation, Ref and case introduction give
\[
 \vdash G_0A,\qquad G_{Sv}A\dashv\vdash A,\qquad
 A\dashv\vdash\D_{w,y}G_yA.
\]
For reverse padding, open D, supply its true zero row, and specialize at 1.

The selector between formulas is
\[
 M_v(s=t,r=q):=(C(v,s,r)=C(v,t,q))
\]
and, with aligned fresh binders,
\[
\begin{aligned}
 M_v(\D_{x,y}A,\D_{x,y}B)&:=\D_{x,y}M_v(A,B),\\
 M_v(E,\D_{x,y}B)&:=\D_{x,y}M_v(G_yE,B),\\
 M_v(\D_{x,y}A,E)&:=\D_{x,y}M_v(A,G_yE),
\end{aligned}
\]
where $E$ is an equation. Maximum input binder depth decreases.
Padding and binder congruence yield
$M_0(A,B)\dashv\vdash A$ and $M_{Sv}(A,B)\dashv\vdash B$.

Put $\delta(v)=C(v,1,0)$ and $p(z)=I(z,0;u.\delta(u))$;
$p(1)=1$ and $p(2)=0$ are computations. With fresh coordinates and unused
$w$, set
\[
\begin{aligned}
 \bot_*&:=(0=1),\\
 A\Rightarrow_*B&:=\D_{w,z}M_z(A,B),\\
 A\wedge_*B&:=\D_{w,z}G_zM_{p(z)}(A,B),\\
 \exists_*x\,A(x)&:=\D_{x,y}G_yA(x),\\
 A\vee_*B&:=\exists_*v\,M_v(A,B),\\
 \forall_*x\,A(x)&:=A(0)\wedge_*\D_{w,z}G_zA(z).
\end{aligned}                                                   \tag{Logic}
\]

Implication introduction is DI$_c$ using its zero row as the assumption;
elimination opens and specializes at 1. Existential introduction supplies
the witness, and elimination opens it freshly and removes the gate.
Disjunction injections use 0 and 1; elimination uses numerical cases on
the selector and the two retained branch implications.
Conjunction introduction uses selector cases and elimination uses rows
1 and 2. The universal keeps zero as a conjunct and supplies successors
in the other conjunct; coverage yields every instance.
Its introduction uses zero and successor substitutions of a proof at a
fresh parameter.

Coverage $(n=0)\vee_*\exists_*m(n=Sm)$ follows by cases and Ref.
From $0=1$, congruence in $v\mapsto C(v,s,t)$ gives every equation;
formula-structural induction gives explosion. These supply all ordinary
intuitionistic rules with the usual eigenvariable conditions.
No excluded middle for arbitrary formulas is added. The ordinary
realization has the explicit falsity bridge $(0=S0)\to\bot$.

\subsection{Scalar uniqueness and the noncircular induction bootstrap}
In U choose $a=f(0)$, use Ref for the base and reverse its conclusion:
\[
 \frac{\Gamma\vdash f(Sz)=h(f(z))}
      {\Gamma\vdash f(n)=I(n,f(0);u.h(u))}.                \tag{UR}\label{eq:ur}
\]
This is exactly the existing scalar uniqueness principle \cite{ur}.
Its equality and logical background has already been recovered.

Appendix~\ref{app:ur} gives the complete inherited proof that UR yields
equation induction. Its dependency order is important:
\begin{enumerate}
\item UR gives $I(n,0;u.Su)=n$ and retraction rigidity. If
$e(0)=0$, $e(e(x))=e(x)$ and $e(S(e(x)))=S(e(x))$, take
$h(u)=e(S(e(u)))$, $r(n)=I(n,0;h)$. Common-update uniqueness gives
$e(r(n))=r(n)$; then $r$ has successor update, so $r(n)=n$ and $e(n)=n$.
For $q(n)=0$, coverage and separation decide the property and the term
$e(n)=C(q(n),n,0)$ provides the retraction split. This derives contextual
induction for zero equations.
\item A scalar predecessor is constructed from finite clocks, a ramp, a
comparison iterator and bounded growth. Common-update uniqueness proves
$P0=0$ and $P(Sn)=n$, without using pairing, a general recursion compiler
or equation induction.
\item Define $d(a,b)=I(b,a;u.Pu)$ and
$\operatorname{sg}(a)=C(a,0,1)$. A scalar diagonal-descent argument and the
explicit accumulator
\[
 \mu(n,a)=d(n,d(n,a)),\quad
 \phi(n,a,b)=\mu(n,a)+\mu(d(n,a),d(b,a))
\]
prove $a+d(b,a)=b+d(a,b)$. Thus
\[
 \Delta(a,b)=\operatorname{sg}(d(a,b)+d(b,a))
\]
is Boolean and $\Delta(a,b)=0\leftrightarrow a=b$ on the entire domain.
\item Translate an equation-induction base and step for $s(n)=t(n)$ to
$\Delta(s(n),t(n))=0$, use zero-equation induction, and reflect back.
The retained context may contain arbitrary first-order formulas.
\end{enumerate}
The substantial scalar construction is inherited and attributed, not
claimed as a new result of U. Its full equations and proofs appear in
the appendix. Once equation induction is recovered, from
$\Gamma,E(z)\vdash E(Sz)$ work under $E(0)$ to derive $E(Sw)$ and use
conditional Q introduction. This recovers the old equation-progress rule.

\section{Equiconsistency and both interpretation composites}
\begin{theorem}[Exact native equivalence]
For every finite sequent of the repository's native calculus $\mathcal N$,
\[
 \Gamma\vdash_{\mathcal N}A
 \quad\Longleftrightarrow\quad
 \nf(\Gamma)\vdash_{\mathcal U}\nf(A).
\]
\end{theorem}
\begin{proof}
Translate native rules to normal forms. Computation becomes Ref,
transport is Section~\ref{sec:equality}, progress is
Section~\ref{sec:bridge}, and the other rules are unchanged.

Conversely, the native calculus proves normalization equalities and
derives U by equation induction on $I(z,a;h)=f(z)$.
Its base uses $f(0)=a$ and symmetry. Its successor step uses the
induction equation, congruence under $h$, computation and the retained
recurrence $f(Sz)=h(f(z))$. All the remaining U-calculus rules are native
rules or derived reflexivity.
\end{proof}

\begin{theorem}[Full first-order iPRA]
The ordinary realization of $\mathcal U$ is equiconsistent and
bi-interpretable with full first-order iPRA. Both numerical domains are
whole, equality is actual, and the two numerical composite isomorphisms
have graph $x=y$.
\end{theorem}
\begin{proof}
Compose the exact native equivalence with the scalar bridge
\cite{bridge}. In iPRA interpret each fixed iterator by primitive
recursion, and Q,D by \eqref{eq:read}. Equation induction validates U.

Conversely, predecessor and equality reflection supply successor
injectivity and characteristics for every fixed quantifier-free formula.
Finite truth tables recover quantifier-free induction. The inherited
scalar construction supplies pairing $\pi$ and projections $L,H$.
For a fixed primitive recursion
\[
 f(0,\vec p)=g(\vec p),\qquad
 f(Sn,\vec p)=h(n,f(n,\vec p),\vec p)
\]
use
\[
\begin{split}
 V(n,\vec p)&=I\bigl(n,\pi(0,g^*(\vec p));
        u.\pi(SL(u),h^*(L(u),H(u),\vec p))\bigr),\\
 f^*(n,\vec p)&=H(V(n,\vec p)).
\end{split}
\]
Equation induction proves the defining equations. This is an external
compiler for fixed descriptions, not an object-language uniform evaluator.

The domain formula is $x=x$ in each direction and equality is unchanged.
Every double-translated operation satisfies its original defining
recurrence. Internal uniqueness identifies the two operations, making
$x=y$ internally total, bijective and operation preserving.

For formulas the logical macros have their ordinary round trips.
The essential paired composite is
\[
 \D_{x,y}B\dashv\vdash
 \exists_*x\bigl(B(x,0)\Rightarrow_*\forall_*y\,B(x,y)\bigr).
\]
Forwards, open D, assume the base, use QE at a fresh parameter and
introduce the macros. Backwards, eliminate the existential and use
the implication and universal instances in DI$_c$.
The Q composite is its conditional universal. Structural induction
supplies both composites for all formulas, including quantified ones.
The explicit falsity bridge preserves contradiction in both directions,
giving equiconsistency.
\end{proof}
This is a written proof using the stated inherited scalar theorems.
It is stronger than mutual interpretability or agreement on standard
numerals, but is not a new proof-assistant verification.


% END INCLUDED


% BEGIN INCLUDED sections/implementation.tex
\section{The separate 502-byte refinement}
This complete program checks the existing guarded public calculus,
not $\mathcal U$:
\begin{lstlisting}[language=Python]
import json

def c(a, g=[0], k=0):
    i = lambda a, t=c, d=0: [i(b, t, d + 5 - len(a)) for b in a] if a * 0 != 0 else a if k <= a < d else d % 2 and i(t, d=3)
    if i(a) in g:
        return
    m, s, u, t = q()
    h = c(m, g, k) or m * 3 or (u and [[s and [s], u, t], i(t, [s, u, t], 3) if s else u]) or (a := i(*a, t, 2))
    c(i((a, i(t, s))[(m, a) == (s and [s, u] or [t] * 3, i(t, u))], [(k := (~-k))]), g + [i(h, (m or 3 % len(a)) and k)], k)
q = json.load(open(0)).pop
c(q())
print('ACCEPT')
\end{lstlisting}
The JSON list is consumed from the end: the first item popped is the
goal, followed by four-field instructions. Acceptance is normal exit
and exactly \file{ACCEPT} followed by a newline. Exceptions reject;
a well-typed unused prefix is allowed.

\subsection{Eager selection and delayed evaluation}
At the old scope floor, let $C_0$ denote the paired elimination test and
$R=i(t,s)$. The old selection
$C_0\mathbin{\mathrm{and}}R\mathbin{\mathrm{or}}a$ becomes $(a,R)[C_0]$.
This Python rewrite is not valid in general. Here the pure opener $i$
preserves a list's root length and sends an integer at root depth zero
either to itself or to zero independently of the image.
When both openings terminate, false $R$ implies $i(t,u)=R$.
If $C_0$ is true, $a=i(t,u)=R$; therefore the old fallback and new
selection agree even then. If the test is false both select $a$.

Extra eager evaluation can reject invalid unused fields, but the
eager-only transformation adds no successful integer/list trace.
Canonical completeness survives because intended source terms are well
scoped and unused fields can be standardized to zero.

The delayed placement is inherited from the 506-byte control.
For a nonzero major, construction of the retained formula does not
depend on the goal. The recursive major check does not mutate its
caller's local scope or context. For a false major, the reconstructed
elimination major is nonempty, so the test is false both before and
after introduction. The final expression computes the selected goal at
the old floor, decrements the floor, and only then performs successor
opening and extends the immutable context.

\subsection{The shorter mode selector}
The selection changes from
\[
 (m\mathbin{\mathrm{or}}\mathord{\sim}\operatorname{len}(a)\bmod2)
   \mathbin{\mathrm{and}}k
 \quad\hbox{to}\quad
 (m\mathbin{\mathrm{or}}3\bmod\operatorname{len}(a))
   \mathbin{\mathrm{and}}k.
\]
At meaningful formula root lengths 1,2,3 both choose $0,k,0$ for a
false major; a true major selects $k$ in both.

For raw aliases, truthiness differs only at odd surplus lengths at least
5. A surplus formula root decodes as administrative truth. If the
branch left the goal unchanged, that call already has a true goal.
If introduction assigned the opened matrix, successful unpacking forces
the entry goal to be an empty truth alias or a D formula with a
surplus-root matrix. The latter is $\D(\mathrm{truth})$, provable with
witness zero. Changed local hypotheses therefore cannot justify a false
entry goal, and they do not escape their call. Empty-root division
errors only remove raw certificates, not canonical proofs.

This is semantic soundness and preservation of canonical proof
compilation, not equality of the entire raw certificate languages.
The strict companion counts exact integer/list and public-goal
validation within its 1172 bytes. The detailed argument, intermediate
controls and provenance are in \file{implementation/compact/}.

\section{Executable evidence and reproduction}
The strict uniqueness checker accepts exactly the record keys
\file{scope}, \file{context}, \file{goal}, \file{proof}.
A closed theorem has scope 0 and context \file{[]}.
Terms are natural-number de Bruijn indices or
\file{["z"]}, \file{["s",t]}, \file{["i",n,a,h]};
the last child binds index 0.
Formulas are \file{["=",s,t]}, \file{["q",B]}, \file{["d",B]}.
D binds $y$ at 0 and $x$ at 1.

The six tags are \file{hyp}, \file{refl}, \file{qe}, \file{de},
\file{di}, \file{u}. U is encoded as
\begin{lstlisting}
["u", f, a, h, n, base_proof, step_proof]
\end{lstlisting}
Both templates have one bound variable. The checker introduces the
recurrence eigenvariable itself and raises passive indices, preventing
hidden counter dependence. All scopes are checked before normalization,
so even ill-scoped subterms erased by computation reject.

Numerical cases in generated certificates expand into DI$_c$, DE, QE
and Hyp. There is no hidden seventh primitive. Eight positive
certificates include symmetry, nontrivial template congruence,
common-source joining using both assumptions, transitivity, and
nondefinitional identity and successor iterator laws.
Eleven targeted invalid cases reject wrong recurrence and base premises,
wrong closure bases, witness escape, unbound variables, Boolean index
aliases and malformed instruction lengths. A separate same-run review
found no concrete defect in the binder handling.

The compact reproduction ran 267 checker processes, with ordinary and
optimized Python, 21 inherited positive/adversarial fixtures, strict
malformed-prefix checks and an intentionally old-only eager certificate.
There were zero unexpected outcomes. These executions validate concrete
implementation behavior; they do not formalize the universal theorems.

From the archive root:
\begin{lstlisting}[language=bash]
python3 reproduce.py
python3 implementation/uniqueness/uniqueness_checker.py \
  < implementation/uniqueness/certificates/common_source_join.json
python3 implementation/compact/strict1172.py \
  < implementation/compact/fixtures/computation_zero_local.json
python3 reference/research_census/reproduce_audit.py
python3 build.py
python3 verify_manifest.py
\end{lstlisting}
The census parses frozen sources without executing them. Checkers and
validation use only the Python standard library; the PDF build needs
LaTeX. Finite recursion or memory limits may reject large valid proofs.
Mathematical completeness uses unbounded-resource semantics; resource
exhaustion is never acceptance.

\section{Limits and further research questions}
The new six-rule basis has a larger checker. The 502-byte program improves
one coordinate for the inherited basis. These two results must not be
combined into a claim that the new mathematical basis achieves that
byte record.

\begin{enumerate}
\item Can seeded crossed uniqueness receive a compact codec that retains
all freshness and counter-independence checks? One complete program below
all current coordinates would meet the strongest form of the request.
\item Can a clause calculus fuse U with witness discharge without
complicating the equality derivation? Existing three-family calculi
already set a demanding comparison.
\item Can the scalar uniqueness bootstrap and this composition be
formalized in Lean? The short elementary equality theorem is a useful
first target.
\item Can normalization share the arity traversal while respecting
sorted positions, protected dummy slots and reindexing? A naive rewrite
of all ternary nodes is invalid.
\item Can a changed-domain relational bi-interpretation reduce the entire
checker? Fixed-language term-definability obstructions do not rule it out;
both definable composite isomorphisms still require proofs.
\end{enumerate}
A smaller checker gives an upper bound on attainable size. An absolute
lower bound requires a fixed implementation language, input contract,
normalization convention and admissible interpretation class.


% END INCLUDED

\appendix
\clearpage
\section{Attributed scalar uniqueness bootstrap}\label{app:ur}
The following sections reproduce the existing scalar uniqueness completion
from the pinned ProveIt sources under their supplied license. They make
the inherited implication used in Section~\ref{sec:bridge} directly
reviewable. Claims of novelty inside the reproduced text refer to that
earlier report. The new elementary equality argument of this report is
Section~\ref{sec:equality}.
\makeatletter
\def\input@path{{reference/scalar_uniqueness/}}
\makeatother

% BEGIN INCLUDED reference/scalar_uniqueness/ur_complete.tex
% Inclusion unit for a complete mathematical report.
% Keep ur_retraction.tex, ur_predecessor.tex and ur_reflection.tex alongside it.
% Requires amsmath, amssymb, amsthm and a theorem environment.

\section{An induction-free scalar presentation}
\label{sec:ur-main}

The new arithmetic principle is unconditional uniqueness of scalar
iteration.  The general idea of replacing induction by recursion
uniqueness is classical.  The result proved here is more specific:
uniqueness for a \emph{single scalar state with no counter argument in
the update} already suffices, in the original scalar language.  It
closes the explicit equality-reflection gap in the incoming
\texttt{scope\_kernel} report.  It does not rely on a general primitive
recursion compiler or on numeric pairing before induction has been
recovered.

\subsection{The precise common background}
Use one numerical sort and fixed finite term syntax
\[
 t ::= x\mid0\mid S(t)\mid I(n,a;u.h).
\]
The binder binds $u$ only in $h$.  Equivalently, each finite earlier-stage
body $h(u,\vec p)$ indexes one first-order function symbol
$I_h(n,a,\vec p)$.  Bodies are not function variables, and the theory
postulates no universal evaluator for primitive-recursive descriptions.
All substitutions avoid capture.

The surrounding logic is intuitionistic first-order logic with actual
equality and genuine falsity $\bot$.  Its usual logical constructors are
$\land,\lor,\to,\forall,\exists,\bot$; negation abbreviates implication
to $\bot$.  In particular, no double-negation elimination or excluded
middle for arbitrary formulas is assumed.

The nonlogical axioms are the universally closed instances
\begin{align*}
 I(0,a;u.h)&=a,\\
 I(Sn,a;u.h)&=h[I(n,a;u.h)/u],\\
 \forall n\,(n=0\lor\exists m\,n=Sm),&\\
 0=S0&\longrightarrow\bot.
\end{align*}
No successor injectivity axiom, equality characteristic, primitive
pairing, or induction axiom occurs in this background.

The uniqueness theory $T_{\rm UR}$ adds the single inference schema
\[
 \boxed{\qquad
 \frac{\Gamma\vdash f(Sz)=h(f(z))}
      {\Gamma\vdash f(t)=I(t,f(0);u.h(u))}.
 \qquad}                                                \tag{UR}
\]
Here $f(z)$ and $h(u)$ are fixed finite term templates.  The update has
no free $z$.  The variable $z$ is fresh for the retained context,
the target $t$, and all passive parameters.  Crucially, the premise is
proved in $\Gamma$ itself: there is no additional induction hypothesis.
The seed $f(0)$ is read from the same template, so there is no separate
base premise.  This is an inference schema over finite syntax, not
quantification over a sort of functions.

Let $T_{\rm eq}$ use the same background and replace UR by contextual
induction on each literal equation $s(n)=r(n)$:
\[
 \frac{\Gamma\vdash s(0)=r(0)\qquad
       \Gamma,s(z)=r(z)\vdash s(Sz)=r(Sz)}
      {\Gamma\vdash s(t)=r(t)},                         \tag{EI}
\]
with the same usual freshness condition.

\begin{theorem}[Scalar uniqueness completion]
For every first-order sequent over the scalar language,
\[
                  \Gamma\vdash_{T_{\rm UR}} A
       \quad\Longleftrightarrow\quad
                  \Gamma\vdash_{T_{\rm eq}} A.
\]
Consequently $T_{\rm UR}$ is bi-interpretable and equiconsistent with full
first-order iPRA.  The arithmetic comparison uses the entire numerical
domain and actual equality; both round trips have identity graph $x=y$.
\end{theorem}

The hard implication is proved below.  The order of the proof prevents
the usual circularity: UR first gives zero-equation induction and a
predecessor, then a scalar accumulator gives equality reflection, and
only then is arbitrary equation induction recovered.  General
primitive recursion is compiled after that point.


% BEGIN INCLUDED reference/scalar_uniqueness/ur_retraction.tex
% Self-contained retraction reduction from scalar UR to zero-equation induction.
\section{Zero-equation induction from uniqueness}
\label{sec:ur-retraction-complete}

The argument in this section is the retraction argument already present in
the incoming \texttt{scope\_kernel} report.  It is reproduced to make the
dependency on equality reflection exact.  It uses neither predecessor nor
a general numerical equality characteristic.

All calculations are under one fixed retained first-order context
$\Gamma$.  Fresh displayed variables are absent from $\Gamma$ and all
passive parameters.  UR first gives
\[
                         I(n,0;u.Su)=n.                  \tag{Id}
\]
Use the template $f(z)=z$ and update $h(u)=Su$; its recurrence premise is
reflexivity.  Common-update uniqueness, obtained by applying UR separately
to two terms and replacing equal seeds, is used below.

\subsection{Retraction rigidity}
Let $e(x)$ be a fixed scalar term template.  Suppose $\Gamma$ proves
\[
 e(0)=0,\qquad e(e(x))=e(x),\qquad e(S(e(x)))=S(e(x)).    \tag{Ret}
\]
Then $\Gamma$ proves $e(t)=t$ for every target $t$.

To see this, put $h(u)=e(S(e(u)))$ and $r(n)=I(n,0;u.h(u))$.
Idempotence in Ret gives $e(h(u))=h(u)$ and $h(e(u))=h(u)$.
Therefore the template $f(n)=e(r(n))$ has initial value zero and
unconditional recurrence
\[
 f(Sn)=e(h(r(n)))=h(r(n))=h(e(r(n)))=h(f(n)).
\]
UR gives $e(r(n))=I(n,0;u.h(u))=r(n)$.
The last equation of Ret now yields
\[
 r(Sn)=h(r(n))=S(e(r(n)))=S(r(n)),\qquad r(0)=0.
\]
A second application of UR, followed by Id, gives $r(n)=n$.
Replacing $r(n)$ by $n$ in the previous equality proves $e(n)=n$.
No induction was used, and no claim that the standard finite successor
orbit exhausts a model was used.

\subsection{A numerical split suffices for induction}
For a fixed formula $A(n)$, suppose
\[
 \Gamma\vdash A(0),\qquad \Gamma,A(z)\vdash A(Sz).
\]
Suppose moreover that $A(n)$ is decidable under $\Gamma$ and that a term
$e(n)$ satisfies $e(n)=n$ under $A(n)$ and $e(n)=0$ under $\neg A(n)$.
Cases give $A(e(n))$: in the first case use $A(n)$, and in the second
use the proved base $A(0)$.  Applying the true-case identity at $e(n)$
gives $e(e(n))=e(n)$.  The induction-step derivation, instantiated at
$e(n)$, gives $A(S(e(n)))$, hence $e(S(e(n)))=S(e(n))$.
The base similarly gives $e(0)=0$.  Retraction rigidity proves $e(n)=n$;
transporting $A(e(n))$ gives $A(n)$.

This is a metatheorem about a justified numerical split.  It supplies no
decision procedure, comprehension principle, or induction schema for
arbitrary quantified formulas.

\subsection{The unconditional zero-equation case}
Write $1=S0$ and $C(v,a,b)=I(v,a;u.b)$.  The two computations give
$C(0,a,b)=a$ and $C(Sv,a,b)=b$.
The separation axiom $0=1\to\bot$ implies $Sv\ne0$: otherwise applying
$w\mapsto C(w,0,1)$ gives $1=0$, hence a contradiction.
Coverage therefore decides $q(n)=0$ for every fixed scalar term $q(n)$.
The term
\[
                             e(n)=C(q(n),n,0)
\]
has exactly the split required in the preceding subsection.
Consequently the following contextual rule is derived:
\[
 \frac{\Gamma\vdash q(0)=0\qquad
       \Gamma,q(z)=0\vdash q(Sz)=0}
      {\Gamma\vdash q(t)=0}.                            \tag{ZI}
\]
Every substitution and case distinction is intuitionistic.  The retained
context can contain arbitrary quantified formulas, provided the fresh
induction variable is absent from it.

% END INCLUDED


% BEGIN INCLUDED reference/scalar_uniqueness/ur_predecessor.tex
% Research fragment for integration by the parent agent.
% Requires amsmath, amssymb, amsthm and theorem/lemma/proof environments.
\section{Scalar iteration uniqueness already defines predecessor}
\label{sec:ur-predecessor}

This section gives the predecessor portion of the completion of the
scalar-uniqueness direction identified in the incoming
\texttt{scope\_kernel} report.  Its new ingredient is a modified scalar
ramp whose invariant is a direct fusion law.  Equality reflection is
proved independently in Section~\ref{sec:ur-reflection-completion}.

Work in ordinary intuitionistic first-order logic with actual equality.
Terms are generated by variables, $0$, $S$, and
$I(n,a;u.h)$, with only $u$ bound in $h$.  Assume the two iterator equations,
\[
 I(0,a;u.h)=a,\qquad I(Sn,a;u.h)=h[I(n,a;u.h)/u],
\]
and numerical coverage $n=0\lor\exists m\,n=Sm$.
Replace equation induction by the contextual scalar uniqueness rule
\[
 \frac{\Gamma\vdash f(Sz)=h[f(z)/u]}
 {\Gamma\vdash f(t)=I(t,f(0);u.h)}\quad(\mathrm{UR}),
\]
where $z$ is fresh for $\Gamma$ and the passive parameters, $f(z)$ is a
term template, and $h(u)$ has no free $z$.  All bodies are finite earlier
syntax.  A useful derived rule says that $f$ and $g$ are equal whenever
they have provably equal initial values and the same unconditional
autonomous update.  Apply UR to both and transport the initial value.
No induction hypothesis is permitted in the recurrence premise.

\begin{theorem}
The preceding theory defines a finite scalar term $P$ and proves
$P(0)=0$ and $P(Sn)=n$.  Consequently it proves successor injectivity.
The derivation does not use separation of $0$ and $S0$.
\end{theorem}

\subsection{Elementary identities by uniqueness}
Write $1=S0$, set $C(n,a,b)=I(n,a;u.b)$, $Z(n)=C(n,1,0)$, and
$a+b=I(b,a;u.Su)$.  Coverage gives $Z(n)=0\lor Z(n)=1$ and
\[
 C(0,a,b)=a,\quad C(Sn,a,b)=b,\quad
 C(e,Sa,Sb)=S(C(e,a,b)).
\]
The last identity follows directly by cases on $e$.
All commutative-monoid laws for addition follow from common-update
uniqueness.  In detail, $n$ has initial value $0$ and update $S$, so
$0+n=n$; $S(a+n)$ has initial value $Sa$ and update $S$, so
$Sa+n=S(a+n)$.  Thus $n+a$ has initial value $a$ and update $S$,
proving $n+a=a+n$.  Finally $a+(b+n)$ has initial value $a+b$ and
update $S$, proving associativity.  These facts justify the usual finite
sum notation $2a=a+a$ and $3a=a+a+a$.

For any fixed body $h$, the following are also common-update
uniqueness proofs:
\[
 I(m+n,a;h)=I(n,I(m,a;h);h),\qquad
 I(Sn,a;h)=I(n,h(a);h).
\]
For the first compare the two sides as functions of $n$, with update $h$
and initial value $I(m,a;h)$; for the second the initial value is $h(a)$.
A constant function fixed by $h$ is likewise the corresponding iterator.

\subsection{Two finite clocks}
Define
\[
 N(n)=I(n,0;u.Z(u)),\qquad j(u)=Z(u)+2N(u),\qquad
 M(n)=I(n,0;u.j(u)).
\]
Coverage on the counter of $N$ proves that $N(n)$ is $0$ or $1$.
Hence $Z(Z(N(n)))=N(n)$ and $N(n)+Z(N(n))=1$ by two finite cases.
The computation rule gives $N(n+2)=N(n)$.
The function $N(2n)$ has initial value $0$ and identity update, whereas
$N(3n)$ has initial value $0$ and update $Z$.  UR therefore gives
\[
 N(2n)=0,\qquad N(2n+1)=1,\qquad N(3n)=N(n).
\]

The range of $j$ lies in $\{0,1,2\}$: at $u=0$ its value is $1$;
at $u=Sv$ the term $Z(u)$ is zero and $2N(u)$ is $0$ or $2$.
Coverage on the counter of $M$ therefore proves that its range also lies
in $\{0,1,2\}$.  Direct computation gives
\[
 j(0)=1,\qquad j(1)=2,\qquad j(2)=0.
\]
Thus $M(n+3)=M(n)$ by three computations and the three cases for
$M(n)$.  UR with the identity update gives $M(3n)=0$, and one further
computation gives $M(3n+1)=1$.

\subsection{The modified ramp and its fusion law}
Use the update
\[
 g(u)=u+N(M(u))+1,\qquad G(n)=I(n,0;u.g(u)).
\]
It differs from the older ramp update $u+M(u)+1$ away from its intended
orbit.  That difference makes the following law valid at every argument:
\[
                         M(g(u))=Z(M(u)).                 \tag{1}
\]
Indeed, when $M(u)=0$, $g(u)=u+1$ and the new $M$ value is $1$;
when $M(u)=1$, $g(u)=u+2$ and the new value is $0$; when $M(u)=2$,
$g(u)=u+1$ and the new value is again $0$.
Consequently $M(G(n))$ has initial value $0$ and update $Z$.
Uniqueness, without any induction hypothesis, proves
\[
                             M(G(n))=N(n).                \tag{2}
\]
As $N$ fixes the two Boolean values, this gives
\[
 G(Sn)=G(n)+N(n)+1,\qquad G(n+2)=G(n)+3,
 \qquad G(n+3)=G(n)+N(n)+4.                               \tag{3}
\]
The last two equations follow by two or three computations and the
two cases for $N(n)$.

The functions $G(2n)$ and $3n$ both start at zero and have update
$v\mapsto v+3$, so UR yields $G(2n)=3n$.
For the less immediate scaling equation, put
\[
 f(n)=G(3n),\qquad k(n)=3G(n)+N(n),\qquad H_0(v)=v+M(v)+4.
\]
Equation (2) and parity scaling give $M(f(n))=N(n)$.
The identities $M(3a)=0$ and $M(3a+1)=1$, together with the two cases
for $N(n)$, give $M(k(n))=N(n)$.
Equations (3) show $f(Sn)=H_0(f(n))$.  In those same two Boolean cases,
\[
 \begin{split}
 k(Sn)&=3G(n)+3N(n)+3+Z(N(n))\\
      &=3G(n)+2N(n)+4=H_0(k(n)).
 \end{split}
\]
Both start at zero.  A second common-update uniqueness proof therefore
gives
\[
                  G(2n)=3n,\qquad G(3n)=3G(n)+N(n).       \tag{4}
\]

\subsection{The scalar comparison recurrence}
Set $H(u)=C(N(u),G(u),u)$.  From (4) and the parity identities,
\[
                 H(2u)=3u,\quad H(1)=1,\quad H(3u)=3H(u).
\]
For the last identity split $N(u)=0$ or $1$; in the zero case use
$G(3u)=3G(u)$, and in the one case both sides retain the corresponding
odd argument.  No division theorem is used.
The functions $I(y,3a;H)$ and $3I(y,a;H)$ have the same initial
value $3a$ and update $H$, hence
\[
 I(y,3a;H)=3I(y,a;H),\qquad I(y,1;H)=1.
\]
Define
\[
 E(x)=I(x,1;u.2u),\qquad L(x,y)=N(I(y,E(x);H)).
\]
The rotation identity and the displayed equations yield
\[
 L(0,y)=1,\qquad L(Sx,0)=0,\qquad L(Sx,Sy)=L(x,y).        \tag{5}
\]
For the final equality, rotate the first iterator step at the initializer
$2E(x)$, replace $H(2E(x))$ by $3E(x)$, move the factor three through
the iterator, and use $N(3v)=N(v)$.

\subsection{Saturated iteration and predecessor}
Define
\[
 B_m(u)=C(L(Su,m),u,Su),\qquad W(t,m)=I(t,0;u.B_m(u)).
\]
Equation (5) and selector compatibility with successor give
\[
 B_{Sm}(Su)=S B_m(u),\qquad B_{Sm}(0)=1.
\]
As functions of $t$, both $W(St,Sm)$ and $S W(t,m)$ start at one
and have update $B_{Sm}$.  Therefore UR proves
\[
                             W(St,Sm)=S W(t,m).           \tag{6}
\]
The function $n\mapsto L(Sn,n)$ starts at zero and has identity update
by (5), so it is identically zero by UR.  By the already universal
equation (6), $n\mapsto W(n,n)$ starts at zero and has update $S$,
so it equals $n$.  It follows that
\[
 W(Sn,n)=B_n(W(n,n))=C(L(Sn,n),n,Sn)=n.                   \tag{7}
\]
Finally put
\[
 P(m)=I\bigl(m,0;u.C(L(SSu,m),u,Su)\bigr).
\]
The zero equation is immediate.  For fixed $m$, the function
\[
 V(t,m)=I\bigl(t,0;u.C(L(SSu,Sm),u,Su)\bigr)
\]
has the same initial value and update as $W(t,m)$, because (5)
identifies its update pointwise with $B_m$.  UR yields $V(t,m)=W(t,m)$.
Using (7),
\[
 P(Sm)=V(Sm,m)=W(Sm,m)=m.
\]
Applying $P$ to $Sx=Sy$ now gives $x=y$.

\subsection{The exact dependency boundary}
Every appeal to induction in the predecessor bootstrap has been replaced
by a specified unconditional common-update uniqueness instance.  All case
distinctions use coverage and finite ranges; no separation axiom, equality
characteristic, or induction on an equation has been smuggled in.
One also obtains $d(x,y)=I(y,x;u.Pu)$ and
\[
 d(Sx,Sy)=d(x,y),\qquad d(a+n,n)=a,
\]
again by common updates $P$ and the identity, respectively.
Consequently addition is cancellative, and
$m(x,y,z)=d(x+z,y)$ is a Mal'cev operation:
$m(x,x,z)=z$ and $m(x,y,y)=x$.

The equality-reflection step needs a different argument.  The older
proof used
\[
 C\bigl(\operatorname{sg}(d(x,y)),x,y+d(x,y)\bigr)=x,
 \qquad \operatorname{sg}(v)=C(v,0,1).
\]
Its familiar proof invokes the equality at $y$ as an induction hypothesis.
The left side has a counter-dependent reset recurrence, not an immediate
autonomous scalar recurrence.  It is not imported into this bootstrap.
Instead, Section~\ref{sec:ur-reflection-completion} constructs an explicit
scalar accumulator using Boolean diagonal descent, proves a maximum
commutativity identity, and obtains actual equality reflection from it.

% END INCLUDED


% BEGIN INCLUDED reference/scalar_uniqueness/ur_reflection.tex
% Inclusion fragment. Requires amsmath, amssymb, and amsthm.
% New scalar-only completion of the UR equality-reflection bootstrap.
% No equation induction, pairing, general PR compiler, or equality test
% is assumed in this fragment. The predecessor bootstrap is a prerequisite.

\section{Actual equality from scalar iteration uniqueness}
\label{sec:ur-reflection-completion}

Work in intuitionistic first-order logic with actual equality, the scalar
iterator computation equations, numerical zero/successor coverage, and
$0\ne S0$. Assume the contextual scalar uniqueness rule
\[
 \frac{\Gamma\vdash f(Sz)=h(f(z))}
      {\Gamma\vdash f(t)=I(t,f(0);u.h(u))},
 \tag{UR}
\]
where $z$ is fresh for $\Gamma,t$ and all passive parameters. The update
$h$ has no active counter argument. Every displayed function in what
follows is a fixed finite scalar term abbreviation; there are no function
variables and no hidden product-state iteration.

The preceding scalar predecessor construction supplies a term $P$ with
$P0=0$ and $P(Sx)=x$. Addition, defined by
$a+b=I(b,a;u.Su)$, satisfies the commutative-monoid equations by common-update
uniqueness. Define
\[
 d(a,b)=I(b,a;u.Pu),\qquad
 Z(a)=C(a,1,0),\qquad \operatorname{sg}(a)=C(a,0,1),
 \qquad C(v,a,b)=I(v,a;u.b).
\]
The following elementary equations use only UR and computation:
\begin{equation}
\begin{aligned}
 d(a,0)&=a,&d(a,Sb)&=P(d(a,b)),\\
 d(Sa,Sb)&=d(a,b),&d(a,a)&=0,&d(0,a)&=0.
\end{aligned}
 \label{eq:ur-monus-basics}
\end{equation}
For cancellation, the two terms $b\mapsto d(Sa,Sb)$ and $b\mapsto d(a,b)$
have common initial value $a$ and common update $P$. The diagonal
$a\mapsto d(a,a)$ has the identity update and initial value zero. The
constant-zero function and $a\mapsto d(0,a)$ have common update $P$ and
initial value zero. Thus no equation-induction assumption is hidden here.

\subsection{A diagonal descent principle for numerical zero tests}

UR implies the following constant-recurrence principle:
\[
 \Gamma\vdash f(Sz)=f(z)\quad\Longrightarrow\quad
 \Gamma\vdash f(t)=f(0).
 \tag{Const}
\]
Apply UR with identity update to $f$ and to the constant function $f(0)$.

Fix a finite arity $r\geq1$ and a scalar term $q(x_1,\ldots,x_r)$.
Suppose that, with any passive parameters fixed,
\[
 q(Sx_1,\ldots,Sx_r)=q(x_1,\ldots,x_r)
\]
and that $q$ is zero whenever any of its displayed arguments is zero.
Then
\[
                         q(x_1,\ldots,x_r)=0.             \tag{DD}
\]
Indeed, coverage in each argument gives
$q(Px_1,\ldots,Px_r)=q(x_1,\ldots,x_r)$: if any argument is zero,
both sides are zero; otherwise the successor identity applies.
Now put
\[
 f(k)=q(d(x_1,k),\ldots,d(x_r,k)).
\]
The defining equation for $d$ gives $f(Sk)=f(k)$, so Const makes $f$
constant. At $k=x_1$ its first argument is $d(x_1,x_1)=0$; hence
$f(x_1)=0$. At $k=0$ it is the desired $q(x_1,\ldots,x_r)$.
This argument uses scalar common-update uniqueness. It does not quantify over
tuples or invoke a simultaneous-recursion rule.

For Boolean-valued terms $p,q$, use
\[
 \operatorname{or}(p,q)=C(p,q,1),\qquad
 \operatorname{xor}(p,q)=C(p,q,Z(q)).
\]
Here Boolean-valued means that the term is provably either $0$ or $1$.
Finite cases show that $\operatorname{xor}(p,q)=0$ if and only if $p=q$.
This finite Boolean fact is not reflection of arbitrary numerical equality.

\subsection{Two Boolean comparison identities}

First, DD gives
\[
 Z(d(Sn,a))=\operatorname{sg}(d(a,n)).
 \tag{B1}\label{eq:ur-strict-complement}
\]
To verify its DD premises, let $q(n,a)$ be the Boolean xor of the two
displayed sides. Simultaneously replacing $n,a$ by $Sn,Sa$ cancels one
successor in each subtraction, so $q(Sn,Sa)=q(n,a)$. If $a=0$, both
sides are zero. If $n=0$, split $a$ into zero and successor: the former
case again gives two zeros, while $a=Sb$ gives
$d(1,Sb)=d(0,b)=0$ and $\operatorname{sg}(Sb)=1$, so both sides are one.
Thus both zero-coordinate boundaries vanish and DD applies.

Second, DD gives
\[
 Z\bigl(d(d(x,a),d(b,a))\bigr)
 =\operatorname{or}\bigl(Z(d(x,a)),Z(d(x,b))\bigr).
 \tag{B2}\label{eq:ur-comparison-join}
\]
Again use the Boolean xor $q(x,a,b)$ of the two sides.
When all three arguments are successors, simultaneous cancellation
reduces the expression to $q(x,a,b)$.
The three zero-coordinate boundaries are as follows.
If $x=0$, both displayed sides are one.
If $a=0$, the left side is $Z(d(x,b))$, and the right side is
$\operatorname{or}(Z(x),Z(d(x,b)))$; coverage on $x$ makes them equal.
If $b=0$, the left side is $Z(d(x,a))$, and the right side is
$\operatorname{or}(Z(d(x,a)),Z(x))$; coverage on $x$ again proves equality.
In both latter arguments, the $x=0$ branch makes the other Boolean one,
whereas the successor branch makes $Z(x)=0$.
Consequently DD proves the vanishing of the xor, and Boolean cases
give the stated identity.

Coverage and the successor equation for addition also give
\[
 \operatorname{or}(\operatorname{sg}(u),\operatorname{sg}(v))
       =\operatorname{sg}(u+v).
 \tag{B3}\label{eq:ur-sign-or}
\]
For $u=0$ this reduces to the sign of $v$; for successor $u$, both
sides are one.

\subsection{An explicit scalar accumulator}

Define the two scalar terms
\begin{equation}
\begin{aligned}
 \mu(n,a)&=d(n,d(n,a)),\\
 \phi(n,a,b)&=\mu(n,a)+\mu(d(n,a),d(b,a)).
\end{aligned}
 \label{eq:ur-scalar-accumulator}
\end{equation}
These are finite compositions of $d$ and addition. In particular,
$\phi$ is not introduced by a counter-dependent recursion.
We now prove its useful counter-dependent equation from its explicit term.

Coverage on $a$ and the monus equations give
\[
                         P(d(Sn,a))=d(n,a).
 \tag{Pd}
\]
At $a=0$ this is $P(Sn)=n$. At $a=Sa'$ it follows from
$d(Sn,Sa')=d(n,a')$ and $P(d(n,a'))=d(n,Sa')$.
Split the value $r'=d(Sn,a)$ into zero and successor.
If $r'=0$, then $d(n,a)=0$ by Pd, so
$\mu(Sn,a)=Sn$ and $\mu(n,a)=n$.
If $r'=Sr$, then $d(n,a)=r$ and
$\mu(Sn,a)=d(Sn,Sr)=d(n,r)=\mu(n,a)$.
In both cases,
\[
             \mu(Sn,a)=\mu(n,a)+Z(d(Sn,a)).              \tag{M}
\]

Put $h=d(b,a)$ and repeat the same split of $r'=d(Sn,a)$.
If $r'=0$, the first argument of the second $\mu$ term is zero at both
$n$ and $Sn$, and $\phi(Sn,a,b)=Sn$, $\phi(n,a,b)=n$.
If $r'=Sr$, the first summand is unchanged, while M gives
\[
                 \mu(Sr,h)=\mu(r,h)+Z(d(Sr,h)).
\]
Associativity therefore proves, in both cases,
\[
 \phi(Sn,a,b)=\phi(n,a,b)+Z(d(d(Sn,a),d(b,a))).
\]
Apply B2 at $x=Sn$, then B1 to each Boolean on its right, and finally B3.
The resulting recursion specification is
\[
 \boxed{\quad
 \phi(0,a,b)=0,\qquad
 \phi(Sn,a,b)=\phi(n,a,b)+\operatorname{sg}(d(a,n)+d(b,n)).
 \quad}                                                   \tag{Acc}
\]
All identities leading to Acc used scalar uniqueness, coverage, and
finite case analysis. No equality-reflection or induction theorem has
yet been used.

\subsection{Commutativity of the arithmetic maximum expression}

Put
\[
                              m(a,b)=a+d(b,a).
\]
The name is descriptive only; its order-theoretic properties are not
assumptions. Four zero/successor cases prove
\[
                  m(a,b)=m(Pa,Pb)+\operatorname{sg}(a+b).
 \tag{Red}\label{eq:ur-max-diagonal-reduction}
\]
At $(0,0)$ both sides are zero. At $(0,Sb)$ and $(Sa,0)$ it is the
successor equation for addition. At $(Sa,Sb)$, simultaneous subtraction
cancellation gives
$m(Sa,Sb)=Sa+d(b,a)=S(m(a,b))$, and the sign of the sum is one.

For fixed $a,b$, form the scalar terms
\begin{align*}
 F(n)&=m(d(a,n),d(b,n))+\phi(n,a,b),\\
 G(n)&=m(d(b,n),d(a,n))+\phi(n,a,b).
\end{align*}
Red, Acc, and commutativity of addition give $F(Sn)=F(n)$ and
$G(Sn)=G(n)$. Hence Const gives
\[
                          F(n)=m(a,b),\qquad G(n)=m(b,a).
\]
At $n=b$, the equality $d(b,b)=0$ and the identities
$m(c,0)=c=m(0,c)$ make both terms equal, by these equations, to
$d(a,b)+\phi(b,a,b)$. It follows that
\[
                    \boxed{\ a+d(b,a)=b+d(a,b).\ }       \tag{Max}
\]
This is the common-accumulator argument for maximum commutativity.
The scalar-specific ingredient is the explicit accumulator in
\eqref{eq:ur-scalar-accumulator} and its derivation through Boolean DD;
no prior availability of arbitrary primitive recursion is required.

\subsection{Reflection and recovery of equation induction}

Define
\[
                       \Delta(a,b)=\operatorname{sg}(d(a,b)+d(b,a)).
\]
Coverage and separation imply $\operatorname{sg}(u)=0\leftrightarrow u=0$.
They also imply $u+v=0\leftrightarrow(u=0\land v=0)$: a successor
in either summand would make the sum a successor, contradicting separation.
Thus $\Delta(a,b)=0$ forces $d(a,b)=0=d(b,a)$. Max then yields
\[
               a=a+d(b,a)=b+d(a,b)=b.
\]
Conversely $a=b$ gives both subtractions zero. Therefore
\[
 \boxed{\quad \Delta(a,b)=0\ \longleftrightarrow\ a=b,\qquad
                  \Delta(a,b)=0\ \lor\ \Delta(a,b)=1.\quad}
 \tag{EqRef}
\]
The equality on the right of EqRef is the logic's actual equality on
the whole numerical domain. Nothing restricts $a,b$ to standard numerals
or to the image of an encoding.

Finally, the independently proved UR retraction lemma gives contextual
induction for all zero equations $q(n)=0$. Given an induction premise
for any equation $s(n)=t(n)$, use EqRef to transfer its base and step
to the zero equation $\Delta(s(n),t(n))=0$, apply that zero-induction
rule, and transfer the conclusion back. This recovers contextual
literal-equation induction in arbitrary retained first-order contexts.
The reverse derivation of UR from literal-equation induction is immediate
by induction on $f(n)=I(n,f(0);h)$. Consequently scalar UR and contextual
equation induction are equivalent over the common scalar computation,
coverage, and separation background, once the preceding UR predecessor
bootstrap is included.

% END INCLUDED


\section{Equivalence, bi-interpretation, and the native calculus}
\label{sec:ur-equivalence-finish}

\subsection{The two proof transformations}
The preceding sections explicitly derive EI inside $T_{\rm UR}$.
For any requested instance $s(n)=r(n)$, replace it by
$\Delta(s(n),r(n))=0$ using EqRef, apply ZI to the transferred base
and step, and use EqRef again at the target.  The context $\Gamma$
has remained fixed throughout.  Thus this argument applies to
contextual induction, not only to closed unconditional theorems.

Conversely, in $T_{\rm eq}$ assume the premise of UR and apply EI to
\[
                       f(n)=I(n,f(0);u.h(u)).
\]
The base is the zero computation and symmetry.  At a successor, the
retained recurrence premise, congruence applied to the induction
equation, and successor computation give
\[
 f(Sn)=h(f(n))=h(I(n,f(0);h))=I(Sn,f(0);h).
\]
This derives UR.  Replacing each use of one rule by its displayed finite
derivation proves equality of all derivable sequents.  The transformation
requires only finite manipulation of fixed terms and proofs.

\subsection{What is inherited in the comparison with iPRA}
The supplied project already proves that this ordinary scalar
$T_{\rm eq}$ is bi-interpretable with full first-order iPRA.  The pinned
reference files are \file{inherited_ordinary_bridge.tex},
\file{equality_reflection.tex}, and
\file{primitive_recursion.tex} under
\file{Research/certified_zero/reference/}, at repository commit
\texttt{11be5aabc6b07262046e2553419bc6b05271d178}.
The exact point at which that bridge is used here is after EI has been
derived independently from UR.

More explicitly, each fixed scalar iterator has its usual
primitive-recursive interpretation in iPRA.  Equation induction proves
its uniqueness, hence validates UR.  In the other direction, EqRef
and the selector assign a Boolean term to every fixed quantifier-free
formula.  Truth tables turn its induction into equation induction.
The inherited arithmetic construction supplies scalar pairing and
projection terms with their inverse equations.  An ordinary primitive
recursion can then be simulated by iterating the paired state consisting
of its counter and current value.  This last use of pairing is legitimate
because equation induction has already been recovered.

All quantifiers in these translations range over the whole numerical
domain.  Equality translates as actual equality.  The two translations
replace only fixed primitive-recursive operation descriptions by the
proved corresponding terms, and the defining recurrences identify the
round trips.  Their numerical isomorphism is $x=y$.  Since the identity
translation between $T_{\rm UR}$ and $T_{\rm eq}$ preserves every sequent,
composing it with the pinned scalar bridge preserves these same whole-
domain and identity-composite properties.  Mutual translation of a
derivation of $\bot$ gives equiconsistency.

\subsection{Replacing progress in the native presentation}
There is also a direct native corollary.  Retain the original terms and
native formula constructors
\[
 A::=(s=t)\mid Q_y A\mid D_{x,y} A,
\]
and their ordinary readings
\[
 Q_y B=B(0)\to\forall y\,B(y),\qquad
 D_{x,y}B=\exists x\,(B(x,0)\to\forall y\,B(x,y)).
\]
Keep Hyp, the two computation equations, equality transport, QE, DE,
and DI by cases; replace equation-progress introduction by UR.
This gives exactly the native sequents of the original calculus.

The logical reconstruction needed for this statement is independent of
the deleted progress rule.  The pinned \texttt{inherited\_choice\_basis.tex}
derives numerical cases, binder congruence, the ordinary intuitionistic
logical interface, and explosion using DI by cases, QE, DE, computation,
and transport.  With crossed one-term transport, the ordinary equality
rules and equation-template transport are first recovered; transport
through $Q$ and $D$ follows by structural induction on formulas using
the case rules.  These are inductions on finite syntax in the metatheory,
not hidden numerical induction instances.

Thus the ordinary UR derivation above can be translated into native
macros without an appeal to progress.  It recovers EI and therefore the
old progress rule.  Conversely, old equation progress gives EI and
the displayed proof of UR.  This establishes the native equivalence.
It does not by itself imply that a particular Python encoding is smaller;
that claim requires the complete checker and its matched measurement.

\subsection{Scope of the result}
The theorem is a written internal derivation, not a proof-assistant
formalization.  Numerical checks of the displayed terms are useful
implementation checks but do not verify the theorem in nonstandard or
intuitionistic models.  No claim of absolute minimality is made.
The mathematical advance is the completion of the restricted scalar
uniqueness basis using an explicit scalar accumulator, without assuming
pairing, counter-aware recursion, or equality reflection along the way.

% END INCLUDED

\clearpage

% BEGIN INCLUDED sections/references.tex
\begin{thebibliography}{99}
\bibitem{current}
V. Reshetnikov, \emph{Concise presentations of intuitionistic primitive
recursive arithmetic}, ProveIt snapshot
\texttt{6ffc49f55d355b548f7a8e1530570aee5de086b3}.
\url{https://github.com/VladimirReshetnikov/ProveIt/tree/6ffc49f55d355b548f7a8e1530570aee5de086b3/Logic/PrimitiveRecursiveArithmetic/ConciseIPRA}.

\bibitem{choice}
ProveIt, \file{certified_zero/reference/inherited_choice_basis.tex}.
Native rules, induction-independent macros and formula composites.
Included under \file{reference/scalar_bridge/}.

\bibitem{bridge}
ProveIt, \file{inherited_ordinary_bridge.tex}, \file{scalar_predecessor.tex},
\file{equality_reflection.tex}, \file{primitive_recursion.tex},
under \file{Research/certified_zero/reference/}.
Included under \file{reference/scalar_bridge/}.

\bibitem{ur}
ProveIt, \file{Research/scalar_uniqueness/ur_complete.tex},
\file{ur_retraction.tex}, \file{ur_predecessor.tex},
\file{ur_reflection.tex}. Reproduced in Appendix~\ref{app:ur}.
\url{https://github.com/VladimirReshetnikov/ProveIt/tree/6ffc49f55d355b548f7a8e1530570aee5de086b3/Logic/PrimitiveRecursiveArithmetic/ConciseIPRA/Research/scalar_uniqueness}.

\bibitem{normalization}
ProveIt, \file{Research/full_normalization/},
\file{Research/computation_normal/}, and Lean modules
\file{FullNormalization}, \file{NativeFullNormal}.
The overview identifies \file{Tm.full}, \file{Tm.full_subst_congr},
and \file{fnprf_iff_nprf}.

\bibitem{goodstein}
R. L. Goodstein, \emph{Recursive Number Theory}, Chapter V:
Formalisations of Primitive Recursive Arithmetic.
Studies in Logic and the Foundations of Mathematics 20 (1957),
104--118. \url{https://doi.org/10.1016/S0049-237X(08)71976-1}.

\bibitem{tait}
W. W. Tait, \emph{Primitive Recursive Arithmetic and Its Role in the
Foundations of Arithmetic: Historical and Philosophical Reflections}.
Author manuscript describing uniqueness and induction.
\url{https://home.uchicago.edu/~wwtx/PRA2.pdf}.
The search index supplied the historical passage; the direct URL returned
404 during this session. The proof here does not depend on that URL.

\bibitem{bytecontrol}
ProveIt,
\file{Research/deterministic_introduction/implementation/code/compact_optional.py},
Git blob \texttt{a0af7afb236b530eb30c63c5ab788151024271b7}.
Included as \file{implementation/compact/base506.py}.

\bibitem{incoming}
ProveIt incoming archives \file{ipra_frontier_review_2026-11-12.zip},
\file{ipra_smaller_envelope.zip}, \file{ipra_orbit_closure.zip}.
Their separate implementation refinements and frozen census are
attributed in the archived provenance files.
\end{thebibliography}


% END INCLUDED

\end{document}

